\documentclass[aps,preprint]{revtex4}%
\usepackage{amsfonts}
\usepackage{amsmath}
\usepackage{amssymb}
\usepackage{graphicx}%
\setcounter{MaxMatrixCols}{30}
%TCIDATA{OutputFilter=latex2.dll}
%TCIDATA{Version=4.00.0.2321}
%TCIDATA{CSTFile=revtex4.cst}
%TCIDATA{Created=Monday, June 10, 2002 12:05:14}
%TCIDATA{LastRevised=Wednesday, June 19, 2002 09:32:30}
%TCIDATA{<META NAME="GraphicsSave" CONTENT="32">}
%TCIDATA{<META NAME="DocumentShell" CONTENT="Articles\SW\REVTeX 4">}
\newtheorem{theorem}{Theorem}
\newtheorem{acknowledgement}[theorem]{Acknowledgement}
\newtheorem{algorithm}[theorem]{Algorithm}
\newtheorem{axiom}[theorem]{Axiom}
\newtheorem{claim}[theorem]{Claim}
\newtheorem{conclusion}[theorem]{Conclusion}
\newtheorem{condition}[theorem]{Condition}
\newtheorem{conjecture}[theorem]{Conjecture}
\newtheorem{corollary}[theorem]{Corollary}
\newtheorem{criterion}[theorem]{Criterion}
\newtheorem{definition}[theorem]{Definition}
\newtheorem{example}[theorem]{Example}
\newtheorem{exercise}[theorem]{Exercise}
\newtheorem{lemma}[theorem]{Lemma}
\newtheorem{notation}[theorem]{Notation}
\newtheorem{problem}[theorem]{Problem}
\newtheorem{proposition}[theorem]{Proposition}
\newtheorem{remark}[theorem]{Remark}
\newtheorem{solution}[theorem]{Solution}
\newtheorem{summary}[theorem]{Summary}
\newenvironment{proof}[1][Proof]{\noindent\textbf{#1.} }{\ \rule{0.5em}{0.5em}}
\begin{document}
\preprint{State Generation}
\title[State Generation]{The Generation of the Convex Set of Quantum States from a Reference State by a
Commutative Semi Group }
\author{B. Weiner}
\affiliation{My Institution}
\keywords{}
\pacs{}

\begin{abstract}


\end{abstract}
\date{06/10/2002}
\startpage{1}
\maketitle
\tableofcontents


\section{Introduction}

In this article we describe how all states, pure and ensemble can be generated
from a reference state by the action of abelian semi groups, which allows one
to express the expectation value of all observable in any state in terms of
"dressed" observables in the reference state. Utilizing this construction and
considering inclusive sequences of semi groups we develope an approximation
theory in the spirit of coherent state approximations. Our construction is
based on the properties of generalized bases, in particular the bases formed
by the generalized eigenvectors of the coordinate operator. Thus in section 2
we review the properties of these bases, in section 3 we describe the abelian
semi groups associated with these bases, in section 4 we describe the
generation of \emph{all} states from a reference state using these semi groups
and in section 6 an approximation theory based on this construction.

\section{Generalized Bases}

Generalized bases of Hilbert spaces have been a cornerstone of quantum
mechanics since their introduction by Dirac \cite{Dirac}, since their
inception they have been given a  firm mathematical foundation in the context
of Distrubution theory, Group theory and Topological Vector Spaces. In this
section we review some of the salient features of their properties.

\section{Coordinate Representation}

If we consider the coordinate representation, all states (unnormalized) in
$\mathcal{H}^{N}$ can be produced in the following manner from a reference
state, $\left\vert \Psi_{ref}\right\rangle $, that is nonzero $a.e.$
\begin{equation}
\left\vert \Psi\left(  \nu,\omega\right)  \right\rangle =\breve{T}\left(
\nu,\omega\right)  \left\vert \Psi_{ref}\right\rangle =\breve{\nu}%
_{N}e^{i\breve{\omega}_{N}}\left\vert \Psi_{ref}\right\rangle
\end{equation}
where
\begin{align}
\left\langle \mathsf{z}^{N}|\Psi_{ref}\right\rangle  &  =\varrho\left(
\mathsf{z}^{N}\right)  e^{iS\left(  \mathsf{z}^{N}\right)  }\nonumber\\
\breve{\nu}_{N}  &  =\int\nu\left(  \mathsf{z}^{N}\right)  \left\vert
\mathsf{z}^{N}\right\rangle \left\langle \mathsf{z}^{N}\right\vert
d\mathsf{z}^{N}\in\mathcal{B}\left(  \mathcal{H}^{N}\right) \\
\breve{\omega}_{N}  &  =\int\omega\left(  \mathsf{z}^{N}\right)  \left\vert
\mathsf{z}^{N}\right\rangle \left\langle \mathsf{z}^{N}\right\vert
d\mathsf{z}^{N}\in\mathcal{B}\left(  \mathcal{H}^{N}\right)
\end{align}
and
\begin{equation}
\left\vert \mathsf{z}^{N}\right\rangle =\left\vert \mathsf{z}_{1}%
\cdots\mathsf{z}_{N}\right\rangle ;\quad\mathsf{z}_{i}\neq\mathsf{z}_{j}%
,\quad1\leq i<j\leq N
\end{equation}
$\lceil$In this formulation it does not matter if the coordinates are equal or
not as for any $N$-particle operator $\breve{\varsigma}_{N}$
\begin{equation}
\breve{\varsigma}_{N}=\int\varsigma\left(  \mathsf{z}^{N}\right)  \prod_{1\leq
i<j\leq N}\delta^{C}\left(  \mathsf{z}_{i}-\mathsf{z}_{j}\right)  \left\vert
\mathsf{z}^{N}\right\rangle \left\langle \mathsf{z}^{N}\right\vert
d\mathsf{z}^{N}=\int\varsigma\left(  \mathsf{z}^{N}\right)  \left\vert
\mathsf{z}^{N}\right\rangle \left\langle \mathsf{z}^{N}\right\vert
d\mathsf{z}^{N}%
\end{equation}
where $\delta^{C}\left(  \mathsf{z}_{i}-\mathsf{z}_{j}\right)  =1-\delta
\left(  \mathsf{z}_{i}-\mathsf{z}_{j}\right)  $, as $\left\vert \mathsf{z}%
_{1}\cdots\mathsf{z}_{N}\right\rangle =0$ if $\mathsf{z}_{i}\neq\mathsf{z}%
_{j}$ for any $i\neq j$.$\rfloor$ Then clearly
\begin{equation}
I_{N}=\int\left\vert \Psi\left(  \nu,\omega\right)  \right\rangle \left\langle
\Psi\left(  \nu,\omega\right)  \right\vert d\mu\left(  \nu,\omega\right)
=\int\breve{T}\left(  \nu,\omega\right)  \left\vert \Psi_{ref}\right\rangle
\left\langle \Psi_{ref}\right\vert \breve{T}\left(  \nu,\omega\right)
^{^{\dagger}}d\mu\left(  \nu,\omega\right)
\end{equation}
where
\begin{equation}
d\mu\left(  \nu,\omega\right)  =\frac{1}{\left\langle \Psi\left(  \nu
,\omega\right)  |\Psi\left(  \nu,\omega\right)  \right\rangle }d\mu_{r}\left(
\nu\right)  d\mu_{i}\left(  \omega\right)
\end{equation}
\qquad\qquad

The strong closure of the space of diagonal integral operators that map
$\mathcal{H}^{N}\rightarrow\mathcal{H}^{N}$ is the space of integral operators
with symmetric integral kernels $T\left(  \mathsf{z}^{N}\right)  $ $\in
L_{\infty}\left(  \mathsf{Z}^{N}\right)  $ /Walter Thirring, \textit{Quantum
Mechanics of Atoms and Molecules -A course in Mathematical Physics 3, pp
50-51; }Springer Verlag , New York 1981/. The space $L_{\infty}\left(
\mathsf{Z}^{N}\right)  $ is a $w^{\ast}$-algebra. The following extract from
Thirring should be noted:

\begin{quotation}
$L_{\infty}\left(  \mathsf{Z}^{N}\right)  ,$\textit{considered as
multiplication operators on} $L_{2}\left(  \mathsf{Z}^{N}\right)  $ \textit{is
maximally abelian. Every function in }$L_{2}\left(  \mathsf{Z}^{N}\right)  $
\textit{that is nonzero} $a.e.$ \textit{is a cyclic vector. Functions that
vanish on some set} $\left[  \Omega^{N}\right]  ^{C}=\mathsf{Z}^{N}-\Omega
^{N}\subset\mathsf{Z}^{N}$ \textit{form invariant subspaces. Thus} $L_{\infty
}\left(  \mathsf{Z}^{N}\right)  $ \textit{is reducible and not a
factor.}\emph{\ Every function in }$L_{2}\left(  \mathsf{Z}^{N}\right)
$\emph{\ that is non zero }$a.e.$\emph{\ is a cyclic vector.}
\end{quotation}

One can see that $T\left(  \nu,\omega\right)  $ (the integral kernel of
$\breve{T}\left(  \nu,\omega\right)  )$ is a parameterization of a
representation of an abelian $w^{\ast}$-algebra, $\mathfrak{A}$, on
$L_{2}\left(  \mathsf{Z}^{N}\right)  $ by $L_{\infty}\left(  \mathsf{Z}%
^{N}\right)  $. This representation is reducible as $L_{2}\left(
\mathsf{Z}^{N}\right)  $ contains the following invariant subspaces
\begin{equation}
L_{2}\left(  \Omega^{N}\right)  =\left\{  \Psi|\Psi\left(  \mathsf{z}%
^{N}\right)  =0\,\text{if \thinspace}\mathsf{z}^{N}\notin\Omega^{N}%
\,\text{,\thinspace}\int_{\Omega^{N}}d\mathsf{z}^{N}\neq0\,\text{and\thinspace
}\Omega^{N}\subset\mathsf{Z}^{N}\right\}
\end{equation}
The operator $\breve{T}\left(  \nu,\omega\right)  $ produces elements that
have amplitudes that belong to $L_{2}\left(  \Omega^{N}\right)  $ when
Support$\left(  \nu\right)  =$ $\Omega^{N}$.

The set of invertible diagonal operators in $\mathfrak{A}$ form a group,
$\mathfrak{G}_{pos}$, that also has a reducible representation on $L_{\infty
}\left(  \mathsf{Z}^{N}\right)  $, as clearly $L_{2}\left(  \Omega^{N}\right)
$ are also invariant subspaces of $\mathfrak{G}_{pos}$. However \emph{it is
not true} that every function in $L_{2}\left(  \mathsf{Z}^{N}\right)  $\ that
is nonzero $a.e.$\ is a cyclic vector for $\mathfrak{G}_{pos}$.

$\lceil$This observation also holds in the discrete case, consider the abelian
algebra generated by $\left\{
%TCIMACRO{\tsum \limits_{1\leq i\leq\infty}}%
%BeginExpansion
{\textstyle\sum\limits_{1\leq i\leq\infty}}
%EndExpansion
\lambda_{i}\left\vert e_{i}\right\rangle \left\langle e_{i}\right\vert
\right\}  $ acting on a vector $a=%
%TCIMACRO{\tsum \limits_{1\leq i\leq\infty}}%
%BeginExpansion
{\textstyle\sum\limits_{1\leq i\leq\infty}}
%EndExpansion
\alpha_{i}\left\vert e_{i}\right\rangle $ that such that $\alpha_{i}%
\neq0\forall i\rfloor$

From the preceding we have that $L_{\infty}\left(  \mathsf{Z}^{N}\right)  $
acting on an $a.e.$ nonzero amplitude produces all (unnormalized) vector
states, the question thus arrises on whether $L_{\infty}\left(  \mathsf{Z}%
^{N}\right)  $ can produce all unnormalized states by acting on an a
$a.e.$nonzero positive kernel $D_{ref}\left(  \mathsf{z}^{N},\mathsf{z}%
^{\prime N}\right)  $ of an operator $\in\mathcal{B}_{1+}\left(
\mathcal{H}^{N}\right)  $ in the following manner%
\begin{equation}
D\left(  \mathsf{z}^{N},\mathsf{z}^{\prime N}\right)  =T\left(  \mathsf{z}%
^{N}\right)  D_{ref}\left(  \mathsf{z}^{N},\mathsf{z}^{\prime N}\right)
T\left(  \mathsf{z}^{\prime N}\right)  ^{\ast}%
\end{equation}
One can start by examining the diagonal elements to obtain
\begin{align}
D\left(  \mathsf{z}^{N},\mathsf{z}^{N}\right)   &  =\left\vert T\left(
\mathsf{z}^{N}\right)  \right\vert ^{2}D_{ref}\left(  \mathsf{z}%
^{N},\mathsf{z}^{N}\right) \nonumber\\
&  \Rightarrow\left\vert T\left(  \mathsf{z}^{N}\right)  \right\vert =\left\{
%
\begin{array}
[c]{c}%
\left(  \frac{D\left(  \mathsf{z}^{N},\mathsf{z}^{N}\right)  }{D_{ref}\left(
\mathsf{z}^{N},\mathsf{z}^{N}\right)  }\right)  ^{\frac{1}{2}};\quad
D_{ref}\left(  \mathsf{z}^{N},\mathsf{z}^{N}\right)  \neq0\\
0;\qquad\qquad\qquad\quad\quad D_{ref}\left(  \mathsf{z}^{N},\mathsf{z}%
^{N}\right)  =0\text{ }%
\end{array}
\right.  \text{ }%
\end{align}
which is well defined $a.e.$ The off diagonal elements will "determine" the
phase of the transformation $T\left(  \mathsf{z}^{N}\right)  $ $=\left\vert
T\left(  \mathsf{z}^{N}\right)  \right\vert e^{i\theta\left(  \mathsf{z}%
^{N}\right)  }$ by
\begin{align}
D\left(  \mathsf{z}^{N},\mathsf{z}^{\prime N}\right)   &  =\left\vert T\left(
\mathsf{z}^{N}\right)  T\left(  \mathsf{z}^{\prime N}\right)  \right\vert
e^{i\left(  \theta\left(  \mathsf{z}^{N}\right)  -\theta\left(  \mathsf{z}%
^{\prime N}\right)  \right)  }D_{ref}\left(  \mathsf{z}^{N},\mathsf{z}^{\prime
N}\right) \nonumber\\
&  \Rightarrow e^{i\left(  \theta\left(  \mathsf{z}^{N}\right)  -\theta\left(
\mathsf{z}^{\prime N}\right)  \right)  }=\frac{D\left(  \mathsf{z}%
^{N},\mathsf{z}^{\prime N}\right)  }{D_{ref}\left(  \mathsf{z}^{N}%
,\mathsf{z}^{\prime N}\right)  \left\vert T\left(  \mathsf{z}^{N}\right)
T\left(  \mathsf{z}^{\prime N}\right)  \right\vert }%
\end{align}
which is well defined as long as
\begin{equation}
\left\vert T\left(  \mathsf{z}^{N}\right)  T\left(  \mathsf{z}^{\prime
N}\right)  \right\vert =0\Longleftrightarrow D\left(  \mathsf{z}%
^{N},\mathsf{z}^{\prime N}\right)  =0
\end{equation}
which means that
\begin{equation}
D\left(  \mathsf{z}^{\prime N},\mathsf{z}^{\prime N}\right)  =0\text{ or
}D\left(  \mathsf{z}^{N},\mathsf{z}^{N}\right)  =0\Longrightarrow D\left(
\mathsf{z}^{N},\mathsf{z}^{\prime N}\right)  =0
\end{equation}
when $D\left(  \mathsf{z}^{N},\mathsf{z}^{\prime N}\right)  \geq0$. First we
see if is true for proper vectors, in that case
\begin{align}
D\left(  \mathsf{v}^{N},\mathsf{v}^{N}\right)   &  =\left\langle
\mathsf{v}^{N}|\breve{D}\mathsf{v}^{N}\right\rangle =\left\langle
\mathsf{v}^{N}|\breve{Q}^{\dagger}\breve{Q}\mathsf{v}^{N}\right\rangle
=0\nonumber\\
&  \Rightarrow\breve{Q}\left\vert \mathsf{v}^{N}\right\rangle =0 \label{c1}%
\end{align}
and
\begin{align}
D\left(  \mathsf{v}^{\prime N},\mathsf{v}^{\prime N}\right)   &  =\left\langle
\mathsf{v}^{\prime N}|\breve{D}\mathsf{v}^{\prime N}\right\rangle
=\left\langle \mathsf{v}^{\prime N}|\breve{Q}^{\dagger}\breve{Q}%
\mathsf{v}^{\prime N}\right\rangle =0\nonumber\\
&  \Rightarrow\breve{Q}\left\vert \mathsf{v}^{\prime N}\right\rangle =0
\label{c2}%
\end{align}
so either condition eq. $\left(  \ref{c1}\right)  $ or eq. $\left(
\ref{c2}\right)  $ clearly implies that
\begin{equation}
\left\langle \mathsf{v}^{N}|\breve{Q}^{\dagger}\breve{Q}\mathsf{v}^{\prime
N}\right\rangle =D\left(  \mathsf{v}^{N},\mathsf{v}^{\prime N}\right)  =0
\end{equation}
To see if this true for the generalized basis we must see if
\begin{equation}
\left\langle \mathsf{z}^{N}|\breve{Q}^{\dagger}\breve{Q}\mathsf{z}%
^{N}\right\rangle =0\Rightarrow\breve{Q}\left\vert \mathsf{z}^{N}\right\rangle
=0
\end{equation}
is true in a rigged Hilbert space. $\lceil\mathcal{V}_{L}\supset
\mathcal{H\supset}\mathcal{V}_{S};$ $\mathcal{V}_{L}$ and $\mathcal{V}_{S}$
are Banach spaces and $\mathcal{V}_{S}$ is dense in $\mathcal{H}$; $\breve{Q}$
acting on $\mathcal{V}_{L}$ is defined by
\begin{equation}
\left(  \breve{Q}^{\ddagger}\mathsf{z}|\varphi\right)  =\left(  \mathsf{z}%
|\breve{Q}\varphi\right)  \forall\left\vert \varphi\right\rangle
\in\mathcal{V}_{S}%
\end{equation}
and $\breve{Q}^{\ddagger}$ is an unique extension of $\breve{Q}$. Note that
$\breve{Q}\left\vert \mathsf{z}^{N}\right\rangle =0$ means
\begin{equation}
\left(  \mathsf{z}|\breve{Q}\varphi\right)  =0\forall\left\vert \varphi
\right\rangle \in\mathcal{V}_{S}%
\end{equation}
which implies that $\breve{Q}^{\ddagger}\left\vert \mathsf{z}^{N}\right\rangle
$ is the zero vector in $\mathcal{V}_{L}$.$\rfloor$ Thus $\left\langle
\mathsf{z}^{N}|\breve{Q}^{\dagger}\breve{Q}\mathsf{z}^{N}\right\rangle =0$ or
$\left\langle \mathsf{z}^{\prime N}|\breve{Q}^{\dagger}\breve{Q}%
\mathsf{z}^{\prime N}\right\rangle =0\Rightarrow\left\langle \mathsf{z}%
^{\prime N}|\breve{Q}^{\dagger}\breve{Q}\mathsf{z}^{N}\right\rangle =0$. Thus
we have that
\begin{equation}
\left\vert T\left(  \mathsf{z}^{N}\right)  T\left(  \mathsf{z}^{\prime
N}\right)  \right\vert =0\Longleftrightarrow D\left(  \mathsf{z}%
^{N},\mathsf{z}^{\prime N}\right)  =0
\end{equation}
and we can define phase differences as
\begin{equation}
e^{i\left(  \theta\left(  \mathsf{z}^{N}\right)  -\theta\left(  \mathsf{z}%
^{\prime N}\right)  \right)  }=\left\{
\begin{array}
[c]{c}%
\frac{D\left(  \mathsf{z}^{N},\mathsf{z}^{\prime N}\right)  }{D_{ref}\left(
\mathsf{z}^{N},\mathsf{z}^{\prime N}\right)  \left\vert T\left(
\mathsf{z}^{N}\right)  T\left(  \mathsf{z}^{\prime N}\right)  \right\vert
};\left\vert T\left(  \mathsf{z}^{N}\right)  T\left(  \mathsf{z}^{\prime
N}\right)  \right\vert \neq0,D_{ref}\left(  \mathsf{z}^{N},\mathsf{z}^{\prime
N}\right)  \neq0\\
\text{\qquad\qquad}0\qquad\text{\qquad\qquad\qquad otherwise\qquad\qquad
\qquad\qquad\qquad}%
\end{array}
\right.
\end{equation}
which is sufficient to define the transformation $D_{ref}\rightarrow D$
\emph{for any }$D$ as long as $D_{ref}\neq0\,a.e.$

The transformations $\hat{T}$ can either decrease or leave invariant the rank
of $\breve{D}^{N}$, thus only if $\breve{D}_{ref}$ has maximal rank will all
density operators be generated by such transformations ($D_{ref}$ need not be
a density operator kernel associated with a fixed number of particles, all
that is needed is that $D_{ref}$ $\in\mathcal{B}_{1+}\left(  \mathcal{H}%
^{N}\right)  $ with $rank\left(  D_{ref}\right)  =\infty$ and thus could be
the kernel of any density operator, including ones describing indeterminate
number of particles.)

\subsection{Second Quantized form of Diagonal Operators}

The operator
\begin{equation}
\breve{T}_{N}=\int T\left(  \mathsf{z}^{N}\right)  \left\vert \mathsf{z}%
^{N}\right\rangle \left\langle \mathsf{z}^{N}\right\vert d\mathsf{z}^{N}%
\end{equation}
is the restriction of the second quantized operator%
\begin{equation}
\breve{T}=\int T\left(  \mathsf{z}^{N}\right)  \Phi\left(  \mathsf{z}%
^{N}\right)  ^{\dagger}\Phi\left(  \mathsf{z}^{N}\right)  d\mathsf{z}^{N}%
\end{equation}
to $\mathcal{H}^{N}$, where
\begin{equation}
\Phi\left(  \mathsf{z}^{N}\right)  =\Phi\left(  \mathsf{z}_{N}\right)
\cdots\Phi\left(  \mathsf{z}_{1}\right)
\end{equation}
and
\begin{equation}
\Phi\left(  \mathsf{z}\right)  ^{\dagger}\left\vert \phi\right\rangle
=\left\vert \mathsf{z}\right\rangle
\end{equation}
As%
\begin{align}
\Phi\left(  \mathsf{z}^{N}\right)  ^{\dagger}\Phi\left(  \mathsf{z}%
^{N}\right)   &  =\Phi\left(  \mathsf{z}_{1}\right)  ^{\dagger}\cdots
\Phi\left(  \mathsf{z}_{N}\right)  ^{\dagger}\Phi\left(  \mathsf{z}%
_{N}\right)  \cdots\Phi\left(  \mathsf{z}_{1}\right) \nonumber\\
&  =\prod_{1\leq i<j\leq N}\delta^{C}\left(  \mathsf{z}_{i}-\mathsf{z}%
_{j}\right)  \Phi\left(  \mathsf{z}_{1}\right)  ^{\dagger}\Phi\left(
\mathsf{z}_{1}\right)  \cdots\Phi\left(  \mathsf{z}_{N}\right)  ^{\dagger}%
\Phi\left(  \mathsf{z}_{N}\right)
\end{align}
we have that%
\begin{align}
\breve{T}  &  =\int T\left(  \mathsf{z}^{N}\right)  \prod_{1\leq i<j\leq
N}\delta^{C}\left(  \mathsf{z}_{i}-\mathsf{z}_{j}\right)  \Phi\left(
\mathsf{z}_{1}\right)  ^{\dagger}\Phi\left(  \mathsf{z}_{1}\right)  \cdots
\Phi\left(  \mathsf{z}_{N}\right)  ^{\dagger}\Phi\left(  \mathsf{z}%
_{N}\right)  d\mathsf{z}^{N}\nonumber\\
&  =\int\sum_{k=0}^{N}\left\{  T_{k}\left(  \mathsf{z}^{N}\right)  \right\}
\prod_{1\leq i<j\leq N}\delta^{C}\left(  \mathsf{z}_{i}-\mathsf{z}_{j}\right)
\Phi\left(  \mathsf{z}_{1}\right)  ^{\dagger}\Phi\left(  \mathsf{z}%
_{1}\right)  \cdots\Phi\left(  \mathsf{z}_{N}\right)  ^{\dagger}\Phi\left(
\mathsf{z}_{N}\right)  d\mathsf{z}^{N}%
\end{align}
where
\begin{equation}
T_{k}\left(  \mathsf{z}^{N}\right)  =\sum_{1\leq i_{1}<i_{2}\cdots<i_{k}\leq
N}T_{k}\left(  \mathsf{z}_{i_{1}}\mathsf{z}_{i_{2}}\cdots\mathsf{z}_{i_{k}%
}\right)
\end{equation}
and all the functions are symmetric.

\section{Pure States \label{pure}}

All pure states can be generated from a reference pure state by the action of
a commutative semi group of operators

\section{Ensemble States \label{mixed}}


\end{document}